% English translation of checked_pages/esq_119_057.tex; source: cote 119, site 57, PDF page 58, typescript page 55.
% AI-translated by GPT-6 Astra (OpenAI, ChatGPT Pro mode).
\scanpage{119/57}
\mspage{55}

probably very provisional, these points are:

a) Combinatorial construction of the Teichmüller Tower.

b) Description of the automorphism group of the profinite
compactification of this tower, and reflection on a characterization
of $\mathrm{I}\!\Gamma = \mathrm{Gal}(\bar{\mathbb{Q}}/\mathbb{Q})$ as a subgroup of the latter.

c) The ``machine for motives'' $S\ell(2,\mathbb{Z})$ and its variants.

d) The anabelian dictionary and the fundamental conjecture
(which is perhaps not so ``out of reach'' as all that !). Among
the crucial points of this dictionary, I envisage the ``profinite
paradigm'' for the fields $\mathbb{Q}$ (cf. b)), $\mathbb{R}$ and $\mathbb{C}$, for which a plausible
formulation remains to be drawn out, as well as a description of the
inertia subgroups of $\mathrm{I}\!\Gamma$, through which the passage from characteristic
zero to characteristic $p > 0$, and to the absolute ring $\mathbb{Z}$, begins.

e) Fermat's problem.

($^5$) I shall nevertheless mention a more delicate piece of work (apart from
the work on cubic complexes mentioned in passing), on
the combinatorial interpretation of the regular maps associated with
the congruence subgroups of $S\ell(2,\mathbb{Z})$. This work was developed
above all with a view to expressing the ``arithmetic'' action of $\mathrm{I}\!\Gamma = \mathrm{Gal}(\bar{\mathbb{Q}}/\mathbb{Q})$
on these ``\uline{congruence maps}'', which takes place, essentially,
through the cyclotomic character of $\mathrm{I}\!\Gamma$. A starting
point was the combinatorial theory of the ``bi-icosahedron'' developed
in a C4 course from purely geometric motivations,
which (as it subsequently turned out) makes it possible to express conveniently
the action of $\mathrm{I}\!\Gamma$ on the category of icosahedral maps (i.e.
congruence maps of index 5).

($^6$) Let us point out in this connection that the isomorphism classes of
compact tame spaces are \unsure{.}the same\tnote{TN-057-01}{The doubtful dot before ``les'' in ``les mêmes'' is retained before the corresponding ``the same''.} as in the ``piecewise
linear'' theory (which, I recall, is \uline{not} a tame theory).
This is, in a sense, a rehabilitation of the ``Hauptvermutung'',\tnote{TN-057-02}{The German name ``Hauptvermutung'' means ``main conjecture''. It is retained as the name used in the source.}
which is ``false'' only because, for historical reasons
that it would probably be interesting to examine more closely, the

\margin{57}
